

\section{Perfil de salida}

En la presente sección se presenta una síntesis del perfil de egreso elaborado en el periodo 2018–2020 por la Comisión de Malla Curricular de la carrera Bach. y Lic. en Matemática, como parte del proceso de transformación curricular que impulsa la Escuela de Matemática. El contenido completo del perfil aprobado se incluye en los anexos de este documento.

El perfil se ha adaptado para reflejar los cuatro grados incluidos esta nueva propuesta.

%\newpage 

\subsubsection{Bachillerato en matemática, énfasis en matemática pura}


\begin{table}[H]
\centering
\renewcommand{\arraystretch}{1.3}
\begin{tabular}{|c|p{9cm}|}
\hline
\cellcolor[HTML]{00C0F3}\textbf{Código(s)} & \cellcolor[HTML]{00C0F3}\textbf{Enunciado de la dimensión declarativa} \\
\hline
CDP01 & Conoce y comprende los fundamentos matemáticos —incluyendo la lógica, la teoría de números y la geometría euclídea— y los principios formales que sustentan el razonamiento y la argumentación rigurosa. \\
\hline
CDP02 & Conoce el uso de software especializado, los fundamentos de la programación y la computación científica, así como las técnicas básicas de análisis numérico para formular, explorar y resolver problemas matemáticos mediante enfoques computacionales. \\
\hline
CDP03 & Conoce los conceptos, métodos y aplicaciones básicos del cálculo y del álgebra lineal como fundamentos de la formación matemática. \\
\hline
CDP04 & Conoce los conceptos y técnicas de las ecuaciones diferenciales, la teoría de la medida, la probabilidad y los procesos estocásticos, así como del análisis complejo y funcional. \\
\hline
CDP05 & Conoce los fundamentos y resultados básicos de la geometría y la topología. \\
\hline
CDP06 & Conoce las estructuras y resultados fundamentales del álgebra abstracta. 
\\
\hline
CDP07 & Conoce las estructuras y resultados fundamentales de algún área de especialización en matemática pura más allá de las antes expuestas. 
\\
\hline
\end{tabular}
\caption{Tabla de conocimientos declarativos o saber conocer.}
\label{tab:declarativa-egreso-integrada}
\end{table}

%\newpage

\begin{table}[H]
\centering
\renewcommand{\arraystretch}{1.3}
\begin{tabular}{|c|p{9cm}|}
\hline
\cellcolor[HTML]{00C0F3}\textbf{Código(s)}  & \cellcolor[HTML]{00C0F3}\textbf{Enunciado de la dimensión procedimental} \\
\hline
CP01 & Aplica pensamiento operacional, abstracto y formal para analizar, resolver y sintetizar problemas matemáticos. \\
\hline
CP02 & Formula conjeturas, resuelve problemas, interpreta resultados y se comunica de manera coherente en lenguaje matemático formal. \\
\hline
CP03 & Utiliza tecnologías y software científico para explorar, validar y resolver problemas matemáticos.\\
\hline
CP04 & Participa activamente en discusiones científicas y divulga el conocimiento matemático de forma accesible a diferentes audiencias.\\
\hline
\end{tabular}
\caption{Tabla de conocimientos procedimentales o saber hacer.}
\label{tab:procedimental-egreso-integrada}
\end{table}

\begin{table}[H]
\centering
\renewcommand{\arraystretch}{1.3}
\begin{tabular}{|c|p{9cm}|}
\hline
\cellcolor[HTML]{00C0F3}\textbf{Código(s)} & \cellcolor[HTML]{00C0F3}\textbf{Enunciado de la dimensión actitudinal} \\
\hline
CA01 & Manifiesta respeto por la diversidad humana, matemática, científica y cultural, fomentando un ambiente inclusivo y de valoración mutua. \\
\hline
CA02 & Se comunica con claridad, colabora en equipos académicos y profesionales, y mantiene una disposición al autoaprendizaje y al liderazgo.\\
\hline
CA03 & Demuestra conciencia social y participación ciudadana, con una visión sociocultural e integral de la sociedad en la que la ciencia y la matemática ocupan un lugar relevante.\\
\hline
CA04 & Ejerce pensamiento crítico, actúa con integridad ética y evalúa con rigor las prácticas científicas y académicas.\\
\hline
CA05 & Utiliza su conocimiento matemático para resolver problemas académicos y profesionales, adaptándolo a distintos contextos y necesidades.\\
\hline
CA06 & Mantiene un interés permanente por los desarrollos de la investigación científica y matemática, valorando la actualización como parte de su formación profesional.\\
\hline
\end{tabular}
\caption{Tabla de conocimientos actitudinales, o saber ser.}
\label{tab:actitudinal-egreso-integrada}
\end{table}


\subsubsection{Bachillerato en matemática, énfasis en matemática aplicada}



\begin{table}[H]
\centering
\renewcommand{\arraystretch}{1.3}
\begin{tabular}{|c|p{9cm}|}
\hline
\cellcolor[HTML]{00C0F3}\textbf{Código(s)} & \cellcolor[HTML]{00C0F3}\textbf{Enunciado de la dimensión declarativa} \\
\hline
CDA01 & Conoce y comprende los fundamentos matemáticos —incluyendo la lógica, la teoría de números y la geometría euclídea— y los principios formales que sustentan el razonamiento y la argumentación rigurosa. \\
\hline
CDA02 & Conoce el uso de software especializado, los fundamentos de la programación y la computación científica, así como las técnicas básicas de análisis numérico para formular, explorar y resolver problemas matemáticos mediante enfoques computacionales. \\
\hline
CDA03 & Conoce los conceptos, métodos y aplicaciones básicos del cálculo y del álgebra lineal como fundamentos de la formación matemática. \\
\hline
CDA04 & Conoce los conceptos y técnicas de las ecuaciones diferenciales, la teoría de la medida, la probabilidad y los procesos estocásticos, así como del análisis complejo y funcional. \\
\hline
CDA05 & Conoce los principios de la matemática aplicada, por ejemplo por medio de la modelación, la estadística o el análisis de datos, para resolver problemas y tomar decisiones en ámbitos como la industria, la banca y otros sectores.
\\
\hline
CDA06 & Conoce las principales tecnologías de aplicación de la matemática, por ejemplo ejemplo la programación y computación científica para aplicar conceptos teóricos a problemas reales.
\\
\hline
CDA07 & Conoce las estructuras y resultados fundamentales de algún área de especialización en matemática aplicada más allá de las antes expuestas. 
\\
\hline
\end{tabular}
\caption{Tabla de conocimientos declarativos o saber conocer.}
\label{tab:declarativa-egreso-integrada}
\end{table}

%\newpage

\begin{table}[H]
\centering
\renewcommand{\arraystretch}{1.3}
\begin{tabular}{|c|p{9cm}|}
\hline
\cellcolor[HTML]{00C0F3}\textbf{Código(s)}  & \cellcolor[HTML]{00C0F3}\textbf{Enunciado de la dimensión procedimental} \\
\hline
CP01 & Aplica pensamiento operacional, abstracto y formal para analizar, resolver y sintetizar problemas matemáticos. \\
\hline
CP02 & Formula conjeturas, resuelve problemas, interpreta resultados y se comunica de manera coherente en lenguaje matemático formal. \\
\hline
CP03 & Utiliza tecnologías y software científico para explorar, validar y resolver problemas matemáticos.\\
\hline
CP04 & Participa activamente en discusiones científicas y divulga el conocimiento matemático de forma accesible a diferentes audiencias.\\
\hline
\end{tabular}
\caption{Tabla de conocimientos procedimentales o saber hacer.}
\label{tab:procedimental-egreso-integrada}
\end{table}


%\newpage

\begin{table}[H]
\centering
\renewcommand{\arraystretch}{1.3}
\begin{tabular}{|c|p{9cm}|}
\hline
\cellcolor[HTML]{00C0F3}\textbf{Código(s)} & \cellcolor[HTML]{00C0F3}\textbf{Enunciado de la dimensión actitudinal} \\
\hline
CA01 & Manifiesta respeto por la diversidad humana, matemática, científica y cultural, fomentando un ambiente inclusivo y de valoración mutua. \\
\hline
CA02 & Se comunica con claridad, colabora en equipos académicos y profesionales, y mantiene una disposición al autoaprendizaje y al liderazgo.\\
\hline
CA03 & Demuestra conciencia social y participación ciudadana, con una visión sociocultural e integral de la sociedad en la que la ciencia y la matemática ocupan un lugar relevante.\\
\hline
CA04 & Ejerce pensamiento crítico, actúa con integridad ética y evalúa con rigor las prácticas científicas y académicas.\\
\hline
CA05 & Utiliza su conocimiento matemático para resolver problemas académicos y profesionales, adaptándolo a distintos contextos y necesidades.\\
\hline
CA06 & Mantiene un interés permanente por los desarrollos de la investigación científica y matemática, valorando la actualización como parte de su formación profesional.\\
\hline
\end{tabular}
\caption{Tabla de conocimientos actitudinales o saber ser.}
\label{tab:actitudinal-egreso-integrada}
\end{table}


\subsubsection{Licenciatura en matemática, énfasis en matemática pura}


En adición a los conocimientos enunciados para el bachillerato, con la licenciatura se tendrían los siguientes.

\begin{table}[H]
\centering
\renewcommand{\arraystretch}{1.3}
\begin{tabular}{|c|p{9cm}|}
\hline
\cellcolor[HTML]{00C0F3}\textbf{Código(s)} & \cellcolor[HTML]{00C0F3}\textbf{Enunciado de la dimensión declarativa} \\
\hline
CDA08 & Conoce estructuras y resultados de nivel intermedio para algunas áreas medulares de especialización en matemática pura. 
\\
\hline
\end{tabular}
\caption{Tabla de conocimientos declarativos o saber conocer.}
\label{tab:declarativa-egreso-integrada}
\end{table}

%\newpage

\begin{table}[H]
\centering
\renewcommand{\arraystretch}{1.3}
\begin{tabular}{|c|p{9cm}|}
\hline
\cellcolor[HTML]{00C0F3}\textbf{Código(s)}  & \cellcolor[HTML]{00C0F3}\textbf{Enunciado de la dimensión procedimental} \\
\hline
CP05 & Aplica pensamiento operacional, abstracto y formal, con fuentes transversales de varias áreas, para formular, analizar y resolver problemas matemáticos. \\
\hline
\end{tabular}
\caption{Tabla de conocimientos procedimentales o saber hacer.}
\label{tab:procedimental-egreso-integrada}
\end{table}


\begin{table}[H]
\centering
\renewcommand{\arraystretch}{1.3}
\begin{tabular}{|c|p{9cm}|}
\hline
\cellcolor[HTML]{00C0F3}\textbf{Código(s)} & \cellcolor[HTML]{00C0F3}\textbf{Enunciado de la dimensión actitudinal} \\
\hline
CA06 & Muestra iniciativa y autonomía para formular preguntas, explorar problemas matemáticos y profundizar en ellos más allá de los procedimientos conocidos.\\
\hline
\end{tabular}
\caption{Tabla de conocimientos actitudinales o saber ser.}
\label{tab:actitudinal-egreso-integrada}
\end{table}

\subsubsection{Licenciatura en matemática, énfasis en matemática aplicada}

En adición a los conocimientos enunciados para el bachillerato, con la licenciatura se tendrían los siguientes.

\begin{table}[H]
\centering
\renewcommand{\arraystretch}{1.3}
\begin{tabular}{|c|p{9cm}|}
\hline
\cellcolor[HTML]{00C0F3}\textbf{Código(s)} & \cellcolor[HTML]{00C0F3}\textbf{Enunciado de la dimensión declarativa} \\
\hline
CDA08 & Conoce estructuras y resultados de nivel intermedio para algunas áreas medulares de especialización en matemática aplicada. 
\\
\hline
\end{tabular}
\caption{Tabla de conocimientos declarativos o saber conocer.}
\label{tab:declarativa-egreso-integrada}
\end{table}

%\newpage

\begin{table}[H]
\centering
\renewcommand{\arraystretch}{1.3}
\begin{tabular}{|c|p{9cm}|}
\hline
\cellcolor[HTML]{00C0F3}\textbf{Código(s)}  & \cellcolor[HTML]{00C0F3}\textbf{Enunciado de la dimensión procedimental} \\
\hline
CP05 & Aplica pensamiento operacional, abstracto y formal, con fuentes transversales de varias áreas, para formular, analizar y resolver problemas matemáticos. \\
\hline
\end{tabular}
\caption{Tabla de conocimientos procedimentales o saber hacer.}
\label{tab:procedimental-egreso-integrada}
\end{table}


\begin{table}[H]
\centering
\renewcommand{\arraystretch}{1.3}
\begin{tabular}{|c|p{9cm}|}
\hline
\cellcolor[HTML]{00C0F3}\textbf{Código(s)} & \cellcolor[HTML]{00C0F3}\textbf{Enunciado de la dimensión actitudinal} \\
\hline
CA06 & Muestra iniciativa y autonomía para formular preguntas, explorar problemas matemáticos y profundizar en ellos más allá de los procedimientos conocidos.\\
\hline
\end{tabular}
\caption{Tabla de conocimientos actitudinales o saber ser.}
\label{tab:actitudinal-egreso-integrada}
\end{table}
